\chapter{面向机械臂底层控制策略的轻量LLM指令微调与部署}\label{chap:finetune}

\section{引言}\label{sec:finetune-intro}
前述章节完成了基于TQC算法的单任务专家策略训练，以及面向大语言模型微调的多任务专家数据集构建。本章研究如何利用\cref{chap:dataset}构建的专家数据集，使用监督指令微调方法对轻量级语言模型进行指令微调，使其从文本化的观测状态生成连续动作序列，并通过vLLM推理引擎实现低延迟部署，满足闭环控制需求。

本章首先在\cref{sec:sft-principle}中介绍监督指令微调的基本原理与LoRA参数高效微调方法。随后，\cref{sec:vllm-deploy}针对LLM推理延迟对闭环控制的影响，引入基于PagedAttention机制的vLLM推理引擎，并给出面向机械臂控制的部署配置。\cref{sec:llm_control-loop}将上述微调与部署模块整合为完整的闭环控制流程，分析各阶段的时序特性。\cref{sec:llm_experiment-design}与\cref{sec:llm_experiment-results}在6个轻量级基座模型上分别开展单任务微调与多任务混合微调实验，从任务成功率与推理延迟两个维度评估系统的控制性能，并分析多任务训练对模型泛化能力的影响。
\section{大语言模型监督式微调原理}\label{sec:sft-principle}
鉴于预训练语言模型并未针对多维度连续物理动作进行优化，需通过监督指令微调(SFT)对齐任务目标，并在算力约束下采用参数高效微调(PEFT)方法。本研究选用LoRA\cite{hu2022lora}方法来实现PEFT。
\subsection{监督指令微调与自回归动作建模}
在常规自然语言处理中，LLMs通过自回归方式进行因果语言建模。给定输入上下文序列$X=\{x_1, x_2, \ldots, x_n\}$，模型旨在最大化目标序列$Y=\{y_1, y_2, \ldots, y_m\}$的条件概率。LLMs基于海量无标注的语料库中进行预训练，能够对常用语料进行拟合。但预训练阶段习得的语言表征难以直接适配特定下游任务的输入格式与输出约束，以及可能对任务目标产生误判，存在分布错位和目标差异等问题。

SFT通过引入特定任务下标注数据集$\mathcal{D}_{\text{SFT}}$，利用有监督学习的方式对预训练模型进行进一步优化。通过提供明确的$(X, Y)$对示例，使模型学习特定任务的输入输出对应关系，完成从文本化观测到动作序列的映射。在本研究中，$\mathcal{D}_{\text{SFT}}$即\cref{chap:dataset}中经过清洗构建的多任务专家数据集。

在本研究的机械臂控制任务中，基于\cref{sec:dialogue-template}所述的处理方法，输入序列$X$包含了任务要求以及当前观测状态，目标序列$Y$为所需生成的动作。指令微调的训练目标即是通过最小化负对数似然损失：
\begin{align}\label{eq:sft-loss}
  \min_{\Theta}\mathcal{L}_{\text{SFT}} &= -\sum_{t=1}^{m} \log P_{\Theta}(y_t | X, y_{<t}), \quad (X,Y) \in \mathcal{D}_{\text{SFT}}
\end{align}
式中$\Theta$代表语言模型的参数，$y_{<t}$表示时刻$t$之前已生成的标记序列。

这种建模方式将传统的连续动作回归问题转换为了序列预测问题。模型的分词器将输入的文字序列映射为对应的标记集合，语言模型则根据前文状态条件，预测后续动作序列，控制机械臂移动和夹爪开合。而由于包含数值信息，本项任务的难点便在于要求模型在微调过程中需要隐式地学习局部词表内的数值连续性规律，因此要求微调过程需要保留模型对数值的敏感度。。

\subsection{LoRA微调原理}
对数十亿参数量级的语言模型进行全参数微调，不仅会导致高昂的显存与计算资源开销，同时也容易因为过拟合导致灾难性遗忘，损害模型原有的语言能力。为此，本研究采用PEFT方法，仅更新参数的增量部分，在保持预训练语言能力的同时适配动作生成任务。PEFT的目标是仅更新模型参数的增量部分，在\cref{eq:sft-loss}所示的优化目标中使用参数量远小于全参数量的子网络$\Phi \ll \Theta$，将增量表示为$\Delta \Theta = \Delta \Theta (\Phi)$，其优化目标可以表示为：
\begin{align}\label{eq:peft-loss}
  \min_{\Phi} \mathcal{L}_{\text{PEFT}} = -\sum_{t=1}^{m} \log P_{\Theta + \Delta \Theta(\Phi)}(y_t | X, y_{<t})
\end{align}

低秩自适应调节(LoRA)是常用的PEFT方法，其核心思想是将预训练模型$\Theta$中的权重矩阵$W_0 \in \mathbb{R}^{d \times k}$分解为两个低秩矩阵$B \in \mathbb{R}^{d \times r}$和$A \in \mathbb{R}^{r \times k}$的乘积，其中$r$远小于$d$和$k$。在微调过程中，原始权重保持不变，仅更新增量矩阵$A$和$B$，使得微调后的权重矩阵为：
\begin{align}
  W = W_0 + \alpha BA
\end{align}
在初始化阶段，矩阵$A$采用高斯分布初始化，而矩阵$B$被初始化为零矩阵，从而确保训练初始状态$\Delta W = 0$，维持预训练模型输出的稳定性。
为平衡基础模型权重与增量矩阵的贡献，LoRA引入了缩放因子$\frac{\alpha}{r}$，其中$\alpha$为超参数。实际前向计算变为：
\begin{align}
  h = W_0x+\frac{\alpha}{r}BAx
\end{align}

Hu等人\cite{hu2022lora}的实验表明，将LoRA模块注入注意力权重$W_q$和$W_v$即可达到与全注意力矩阵注入相当的效果，且较低的秩$r$足以使模型在下游任务上收敛至较优解。但在后续实验中发现，对于本研究的数值密集型动作生成任务，仅注入注意力层不足以引导模型学习数值密集型的动作生成规律。因此，本研究将LoRA模块的注入范围扩展至注意力层的全部投影矩阵$(W_q, W_k, W_v, W_o)$以及前馈网络层的门控与升降维投影矩阵$(W_{\text{gate}}, W_{\text{up}}, W_{\text{down}})$，以在更广泛的参数空间中引导模型学习动作序列的生成分布。

\begin{figure}[htbp]
\centering
\includegraphics[width=0.70\textwidth]{figures/lora_structure}
\caption{LoRA结构示意图}\label{fig:lora-structure}
\end{figure}

\section{面向底层控制的vLLM低延迟部署机制}\label{sec:vllm-deploy}
在强化学习的单任务专家模型中，MLP前向传播仅需微秒级时间，满足底层控制的高频需求。而引入大语言模型作为通用控制器后，参数规模大幅增加，推理延迟显著升高。本节引入vLLM\cite{kwon2023efficient}推理引擎进行部署优化，以提升推理速度。

\subsection{推理阶段KV缓存与显存瓶颈}
基于Transformer架构的大语言模型在自回归推理阶段，需缓存历史生成的Key和Value张量以避免重复计算。设模型共有$L$层、$h$个注意力头、每头维度为$d_h$，当前序列长度为$s$，则KV缓存的显存占用为：
\begin{align}
  M_{\text{KV}} = 2 \times L \times h \times d_h \times s \times b
\end{align}
式中$b$为每个元素的字节数，因子2分别对应Key和Value。可见，KV缓存的显存需求与序列长度$s$线性相关。在传统推理框架中，通常为每个请求预分配一块与最大序列长度等大的连续显存空间。这种静态分配策略为每个请求预分配最大长度的连续显存，在请求数较多时造成显存浪费，限制了单卡并发处理能力。

\subsection{基于分页注意力的显存管理}
Kwon等人\cite{kwon2023efficient}提出PagedAttention机制，借鉴操作系统虚拟内存的分页管理思想，将KV缓存组织为非连续的固定大小内存块(block)。每个block存储固定数量token的Key和Value向量，不同block在物理显存中无需相邻，通过block表维护逻辑block到物理block的映射关系。

PagedAttention将KV缓存的逻辑序列划分为若干block，每个block包含$B_s$个token的KV数据。对于一个长度为$s$的序列，其逻辑block数为$\lceil s / B_s \rceil$。block表记录了每个逻辑block对应的物理block地址，注意力计算时通过block表索引即可访问非连续物理显存中的KV数据。

该机制实现了显存按需分配，仅在实际生成新token时才分配新的block，减少了预分配的浪费。同时，不同请求的KV缓存可以共享物理block，在多任务并行推理时提升显存利用率。

vLLM还采用连续批处理策略，在迭代级别而非请求级别进行批处理调度。某个请求生成完毕后，其占用的block立即释放并分配给新请求，无需等待同批次其他请求完成，从而减少GPU空闲等待。

\subsection{面向机械臂控制的推理部署配置}
在本研究的机械臂控制场景中，LLM以闭环方式参与每一步决策，推理延迟直接影响控制频率与系统稳定性。为此，本研究基于vLLM引擎进行了针对性的部署配置。

在模型加载阶段，考虑到机械臂控制任务的输入序列由结构化的观测状态描述与任务指令构成，输出为4维动作数值序列，整体序列长度较短，因此将最大模型序列长度$\texttt{max\_model\_len}$设置为1024以限制KV缓存的显存预分配上限。同时，将GPU显存利用率$\texttt{gpu\_memory\_utilization}$设置为0.3，在单卡多任务并存的实验环境下为其他进程预留显存。再启用$\texttt{enforce\_eager}$模式以跳过CUDA Graph捕获阶段，减少首次推理的预热开销。

在采样策略方面，由于机械臂控制要求同一观测状态下的动作输出具有确定性以避免末端执行器抖动，将采样温度$\texttt{temperature}$设置为0.0，即采用贪婪解码策略，在每个时间步选择概率最大的token作为输出。同时，将最大生成token数$\texttt{max\_tokens}$设置为64，以限制动作序列的生成长度并保障推理效率。


\section{大语言模型驱动的闭环控制流程建模}\label{sec:llm_control-loop}
本节将各模块整合为完整的闭环控制流程，并对其时序特性进行分析。

\subsection{闭环控制流程}
本研究构建的LLM驱动闭环控制系统如\cref{fig:control-loop}所示，其完整执行流程由以下五个阶段构成：

\begin{enumerate}
  \item 状态观测：在时间步$t$，仿真环境Gymnasium-Robotics返回物理状态向量$s_t$，包含末端执行器位姿、目标物体状态及当前目标位置等信息。
  \item 文本编码：基于\cref{sec:data-transformation}设计的异构数据转化机制，将$s_t$中的各物理变量映射为具备语义信息的结构化文本描述，并按照\cref{sec:dialogue-template}的对话模板封装为完整提示词$X_t$。
  \item 动作推理：将$X_t$输入经LoRA微调后部署于vLLM引擎的语言模型$\Theta_{\text{LoRA}}$，模型通过自回归方式生成动作标记序列$\hat{Y}_t$。
  \item 动作解析：通过正则匹配从$\hat{Y}_t$中提取形如$[a_1, a_2, a_3, a_4]$的浮点数序列，其中$a_1, a_2, a_3$分别为末端执行器在$x, y, z$方向的增量位移，$a_4$为夹爪开合控制量，解析结果转换为连续动作向量$\hat{a}_t \in \mathbb{R}^4$。若模型输出无法匹配预期格式，则输出零向量作为安全回退动作。
  \item 环境执行：将$\hat{a}_t$输入仿真环境执行，产生下一时间步的状态$s_{t+1}$与奖励信号$r_t$，完成一个闭环周期。
\end{enumerate}

\begin{figure}[htbp]
\centering
\includegraphics[width=0.95\textwidth]{figures/LLM_control_workflow}
\caption{LLM驱动闭环控制流程图}\label{fig:control-loop}
\end{figure}

\subsection{控制时序分析}
上述闭环流程的每一步均引入了不可忽略的计算开销。设状态观测与环境执行的时间分别为$t_{\text{obs}}$与$t_{\text{exec}}$，文本编码时间为$t_{\text{enc}}$，LLM推理时间为$t_{\text{infer}}$，动作解析时间为$t_{\text{parse}}$，则单个控制周期的总延迟为：
\begin{align}\label{eq:control-latency}
  T_{\text{step}} = t_{\text{obs}} + t_{\text{enc}} + t_{\text{infer}} + t_{\text{parse}} + t_{\text{exec}}
\end{align}
其中，$t_{\text{infer}}$占比最大。与\cref{chap:reinforce}中MLP专家模型的微秒级前向传播相比，LLM的推理延迟高出数个数量级不过，本研究面向的机械臂操作任务对控制频率的需求在1\textasciitilde10 Hz范围内，LLM的推理延迟尚可满足。vLLM引擎对$t_{\text{infer}}$的优化效果将在\cref{sec:llm_experiment-design}及\cref{sec:llm_experiment-results}中通过实验量化。

\section{实验设计}\label{sec:llm_experiment-design}
本节评估不同基座模型经LoRA微调后在机械臂操作任务上的控制性能，主要关注成功率与推理延迟。

\subsection{基座模型}
本研究选取6个参数量在0.6B至1.3B之间的轻量级开源大语言模型，基本信息如\cref{tab:base-models}所示。

\begin{table}[htbp]
  \centering
  \caption{实验所用基座模型}\label{tab:base-models}
  \begin{tabular}{lccc}
    \toprule
    模型 & 参数量 & 开发机构 & 架构 \\
    \midrule
    Qwen3-0.6B\cite{yang2025qwen3} & 0.6B & 阿里云 & Decoder-only \\
    Qwen3.5-0.8B\cite{yang2025qwen3} & 0.8B & 阿里云 & Decoder-only \\
    Llama3.2-1B\cite{grattafiori2024llama} & 1B & Meta & Decoder-only \\
    Gemma3-1B\cite{team2025gemma} & 1B & Google & Decoder-only \\
    OLMo2-1B\cite{olmo2024olmo2} & 1B & AI2 & Decoder-only \\
    Phi-1.5\cite{li2023phi} & 1.3B & Microsoft & Decoder-only \\
    \bottomrule
  \end{tabular}
\end{table}

所有模型均采用Decoder-only架构，参数量在1B级别，适合在单张消费级GPU上推理。

\subsection{微调与部署配置}
所有模型统一采用如\cref{tab:training-config}所示的微调配置。LoRA模块的注入范围覆盖注意力层的全部投影矩阵$(W_q, W_k, W_v, W_o)$以及前馈网络层的门控与升降维投影矩阵$(W_{\text{gate}}, W_{\text{up}}, W_{\text{down}})$，秩$r$设为16，缩放因子$\alpha$设为32。

\begin{table}[htbp]
  \centering
  \caption{LoRA微调与训练超参数配置}\label{tab:training-config}
  \begin{tabular}{ll}
    \toprule
    配置项 & 取值 \\
    \midrule
    LoRA秩 $r$ & 16 \\
    LoRA缩放因子 $\alpha$ & 32 \\
    LoRA Dropout & 0.05 \\
    目标矩阵 & $W_q, W_k, W_v, W_o, W_{\text{gate}}, W_{\text{up}}, W_{\text{down}}$ \\
    \midrule
    训练轮数 & 3 \\
    学习率 & $5 \times 10^{-5}$ \\
    批大小 & 4 \\
    梯度累积步数 & 4 \\
    最大序列长度 & 512 \\
    精度 & BF16 \\
    \midrule
    训练硬件 & NVIDIA RTX 4070 Ti Super \\
    \bottomrule
  \end{tabular}
\end{table}

在部署阶段，所有模型均采用\cref{sec:vllm-deploy}所述的vLLM引擎配置，采样温度设为0.0以输出当前最优结果，最大生成token数设为64。

\subsection{评估方案}
性能评估采用在线仿真闭环测试，每个模型在每个任务上独立运行500个回合。实验以多任务混合微调为主，以单任务独立微调作为消融对照。评估指标包括：

\begin{itemize}
  \item 成功率：在500个回合中，目标物体与目标地点距离小于0.05m的回合占比。
  \item 格式解析率：模型输出被正则表达式成功解析为4维浮点数向量的比例，能够反映模型对输出格式约束的符合程度。
  \item 单步推理延迟：从提示词输入vLLM引擎到动作向量解析完成的平均时间，基于500个回合的逐帧记录取均值。
\end{itemize}

作为性能上界参考，\cref{sec:expert-training}中训练的MLP专家模型在同样4个任务上经1000个回合评估的成功率也纳入\cref{tab:experiment-results}进行对比。

\section{实验结果与分析}\label{sec:llm_experiment-results}
基于\cref{sec:llm_experiment-design}所设计的实验，各模型在多任务混合微调与单任务独立微调两种设定下的闭环控制性能如\cref{tab:experiment-results}所示。

\begin{table}[htbp]
  \centering
  \caption{各模型多任务混合微调与单任务独立微调闭环控制性能对比}\label{tab:experiment-results}
  \small
  \begin{tabular}{llccccccc}
    \toprule
    \multirow{2}{*}{模型} & \multirow{2}{*}{\shortstack{微调\\设定}} & \multicolumn{4}{c}{成功率 (\%)} & \multirow{2}{*}{\shortstack{平均\\成功率(\%)}} & \multirow{2}{*}{\shortstack{格式\\解析率(\%)}} & \multirow{2}{*}{\shortstack{平均延迟\\(ms)}} \\
    \cmidrule(lr){3-6}
    & & Reach & Push & PickAndPlace & Slide & & & \\
    \midrule
    MLP专家模型 & — & 100.0 & 97.9 & 99.0 & 90.6 & 96.9 & 100.0 & $<$1 \\
    \midrule
    \multirow{2}{*}{Qwen3.5-0.8B} & 多任务 & 99.4 & 23.6 & 65.4 & 31.2 & 54.9 & 100.0 & 466.6 \\
                                   & 单任务 & 99.4 & 29.6 & 42.0 & 28.6 & 49.9 & 100.0 & 460.8 \\
    \multirow{2}{*}{Gemma3-1B} & 多任务 & 98.6 & 19.2 & 56.6 & 26.4 & 50.2 & 98.0 & 401.2 \\
                                & 单任务 & 90.0 & 18.8 & 20.6 & 5.4 & 33.7 & 100.0 & 415.5 \\
    \multirow{2}{*}{Llama3.2-1B} & 多任务 & 78.8 & 14.2 & 53.4 & 21.4 & 41.95 & 100.0 & 131.2 \\
                                  & 单任务 & 70.2 & 12.8 & 26.8 & 6.6 & 29.1 & 100.0 & 128.9 \\
    \multirow{2}{*}{Qwen3-0.6B} & 多任务 & 99.4 & 13.8 & 30.4 & 19.4 & 40.75 & 100.0 & 189.1 \\
                                 & 单任务 & 98.8 & 10.6 & 44.4 & 25.2 & 44.75 & 100.0 & 181.0 \\
    \multirow{2}{*}{OLMo2-1B} & 多任务 & 37.4 & 9.2 & 16.8 & 7.6 & 17.75 & 100.0 & 138.9 \\
                               & 单任务 & 8.8 & 12.8 & 4.4 & 1.8 & 6.95 & 100.0 & 135.2 \\
    \multirow{2}{*}{Phi-1.5} & 多任务 & 6.4 & 6.4 & 2.4 & 0.4 & 3.9 & 100.0 & 136.4 \\
                              & 单任务 & 1.0 & 7.4 & 1.8 & 0.0 & 2.55 & 100.0 & 133.4 \\
    \bottomrule
  \end{tabular}
\end{table}

\subsection{多任务混合微调性能分析}
如\cref{tab:experiment-results}及\cref{fig:success-rate-comparison}所示，在多任务混合微调设定下，Qwen3.5-0.8B以54.9\%的平均成功率取得了最优性能，在PickAndPlace任务上达到65.4\%，为所有模型中最高。Gemma3-1B以50.2\%的平均成功率位居第二，在PickAndPlace任务上达到56.6\%。Llama3.2-1B与Qwen3-0.6B紧随其后，平均成功率分别为41.95\%与40.75\%。OLMo2-1B的平均成功率为17.75\%，较其单任务基线有显著提升。Phi-1.5的平均成功率仅为3.9\%，表明该基座模型的预训练表征与数值密集型动作生成任务存在较大差距。作为性能上界参考，MLP专家模型的平均成功率达到96.9\%，所有LLM控制策略与之差距仍然较大。

\begin{figure}[htbp]
\centering
\includegraphics[width=0.95\textwidth]{figures/success_rate_comparison}
\caption{各模型多任务混合微调成功率}\label{fig:success-rate-comparison}
\end{figure}

所有模型在Reach任务上的成功率普遍较高，因为Reach仅涉及末端执行器的运动学控制，无需处理物体接触与碰撞。随着任务复杂度增加，成功率普遍下降，其中Slide任务下降最为明显。如\cref{sec:data-quality-analysis}所述，Slide任务对初始接触冲量极为敏感，微小的动作偏差会被动力学积分放大，这一特性同样影响了LLM控制策略的表现。

在格式解析率方面，除Gemma3-1B在多任务设定下的格式解析率为98.0\%外，所有模型在两种设定下的格式解析率均达到或接近100\%，表明监督式微调已使模型充分学习到动作输出的格式约束，能够稳定生成符合$[a_1, a_2, a_3, a_4]$结构的文本序列。

\subsection{单任务消融对比分析}
对比同一模型在两种设定下的性能，多任务训练对不同基座模型的影响各不相同。其中，Gemma3-1B的多任务增益最大，平均成功率从单任务设定的33.7\%提升至多任务设定的50.2\%，提升了16.5个百分点。该提升主要体现在PickAndPlace与Slide两个复杂任务上，分别从20.6\%和5.4\%提升至56.6\%和26.4\%，表明跨任务知识迁移对中等表现模型具有显著的正向效应。

Llama3.2-1B同样展现出显著的多任务增益，平均成功率从单任务的29.1\%提升至多任务的41.95\%，提升了12.85个百分点。该模型在PickAndPlace与Slide任务上分别提升了26.6与14.8个百分点，进一步验证了多任务训练对中等表现模型的正向作用。

在单任务设定中表现最优的Qwen3.5-0.8B在多任务微调中也获得了5.0个百分点的提升，平均成功率从49.9\%提高至54.9\%。该提升主要来自PickAndPlace任务上23.4个百分点的增长（42.0\%$\to$65.4\%），而Push任务则出现了6.0个百分点的性能衰减。这表明即使是高表现模型，也能从多任务训练中获益，尽管部分任务间仍存在参数竞争。

Qwen3-0.6B在多任务微调中出现了4.0个百分点的平均成功率下降，主要体现在PickAndPlace任务从44.4\%降至30.4\%。这种现象揭示了在模型容量有限时，多任务训练中的资源分配竞争可能抑制部分任务的表现。

总体来看，多任务混合微调在大多数模型上有效，其中Gemma3-1B与Llama3.2-1B增益最大，说明单一LLM可以同时掌握多项操作技能。但Qwen3-0.6B的性能下降表明，在模型容量有限时，简单的数据混合策略仍有局限。

\subsection{推理延迟分析}
各模型在多任务混合微调设定下的单步推理延迟对比如\cref{tab:latency-vllm}所示。为评估\cref{sec:vllm-deploy}所述vLLM推理引擎的实际加速效果，本节同时记录了各模型在原生HuggingFace推理框架下的延迟作为对照。

\begin{table}[htbp]
  \centering
  \caption{vLLM推理引擎对各模型推理延迟的加速效果}\label{tab:latency-vllm}
  \small
  \begin{tabular}{lccc}
    \toprule
    模型 & 不使用vLLM (ms) & 使用vLLM (ms) & 加速比 \\
    \midrule
    Qwen3-0.6B & 320.80 & 189.1 & 1.70 \\
    Qwen3.5-0.8B & 466.60 & 338.8 & 1.38 \\
    Llama3.2-1B & 164.70 & 131.2 & 1.26 \\
    Gemma3-1B & 465.55 & 401.2 & 1.16 \\
    OLMo2-1B & 192.49 & 138.9 & 1.39 \\
    Phi-1.5 & 169.17 & 136.4 & 1.24 \\
    \bottomrule
  \end{tabular}
\end{table}

由\cref{tab:latency-vllm}可知，vLLM推理引擎对所有模型均有加速效果。其中，Qwen3-0.6B的加速效果最为显著，推理延迟从320.80ms降至189.1ms，加速比达到1.70倍；Qwen3.5-0.8B的延迟从466.60ms降至338.8ms，加速比为1.38倍。Gemma3-1B的加速效果相对有限，加速比仅为1.16倍，这可能与其模型内部的注意力计算模式对PagedAttention的适配程度有关。

经vLLM加速后，Llama3.2-1B实现了131.2ms的最短推理延迟，对应约7.6 Hz的控制频率。Qwen3-0.6B、OLMo2-1B和Phi-1.5的延迟在136\textasciitilde189ms范围内，对应约5.3\textasciitilde7.3 Hz的控制频率。Qwen3.5-0.8B和Gemma3-1B的延迟分别约为338.8ms和401.2ms，对应约2.5\textasciitilde3.0 Hz的控制频率。

\begin{figure}[htbp]
\centering
\includegraphics[width=0.92\textwidth]{figures/latency_comparison}
\caption{各模型多任务微调设定下推理延迟与对应控制频率对比}\label{fig:latency-comparison}
\end{figure}

上述延迟均在机械臂操作任务1\textasciitilde10 Hz的可接受范围内，满足实时控制需求。但与MLP专家模型的微秒级推理相比，LLM的推理延迟仍高出约5个数量级。这一差距源于通用语言模型与专用控制网络在架构复杂度上的差异。

\subsection{小结}
本章实验在6个轻量级基座模型上评估了LoRA微调方案的控制性能。多任务混合微调设定下，Qwen3.5-0.8B以54.9\%的平均成功率领先，Gemma3-1B与Llama3.2-1B分别获得16.5与12.85个百分点的多任务增益，表明单一LLM可以同时掌握多项操作技能。单任务微调中，Qwen3-0.6B以44.75\%取得最优性能，但多任务微调时下降4.0个百分点，说明模型容量有限时数据混合策略存在局限。所有模型的推理延迟均在1\textasciitilde10 Hz范围内。不过，LLM控制策略与MLP专家模型之间仍有较大性能差距，任务条件化的微调策略或扩大模型规模是可能的改进方向。